\section{Appendix}

\subsection{Numerical Jacobian computation}
The Jacobian $J_x$ is the matrix of all first-order partial derivatives of a vector-valued function.
For the prediction and measurement model, the Jacobians of  the functions $f$ and $h$ (\footnote{The following equations are shown for an ``$f$", but are of course the same for $f$ and $h$}) are needed with respect to the state $x$ (where $n = \mathrm{length}(x)$).
The Jacobian of any function with respect to $x$ is defined as \cite{martins2021}
\begin{align}
    J_x = \frac{\partial f}{\partial x}
    = \begin{bmatrix}
        \frac{\partial f_1}{\partial x_1} & \cdots & \frac{\partial f_1}{\partial x_n} \\
        \vdots & \ddots & \vdots \\
        \frac{\partial f_m}{\partial x_1} & \cdots & \frac{\partial f_m}{\partial x_n}
    \end{bmatrix}
\end{align}
Jacobians may be derived analytically from the function, but this is often cumbersome and prone to errors.
Therefore, it is common practice to compute the matrix numerically.
A frequently used method is using \textit{finite differences}, however this is known to be not very efficient nor accurate, and has issues with numerical instability \cite{martins2021}.

\subsubsection{Complex-step differentiation}
A more efficient way to calculate Jacobian matrices numerically is to use \textit{complex step differentiation}.
In this method with step size $\varepsilon$, the $k$-th column of the matrix ($J_{x_k}$) is \cite{martins2003}
\begin{align}
    J_{x_k} &= \mathrm{Im} \left( \frac{f(t, x + \mathrm{i} \varepsilon e_k, u)}{\varepsilon} \right)
    \label{eq:appendix_jacobian_complex_step}
\end{align}
where $e_k$ is the $k$-th unit vector in the state-space (i.e. $e \in X$, and $e_1 = [\begin{smallmatrix} 1 & 0 & \cdots & 0 \end{smallmatrix}]^T$, $e_2 = [\begin{smallmatrix} 0 & 1 & \cdots & 0 \end{smallmatrix}]^T$, etc.).
This must be computed for each entry of $x$, so that the Jacobian is 
\begin{align}
    J_x &= \begin{bmatrix}
        J_{x_1} & J_{x_2} & \cdots & J_{x_n} 
    \end{bmatrix} 
\end{align}
The evaluation of the functions with a complex state vector $[\begin{smallmatrix} x_1 + \mathrm{i} & x_2 & \cdots & x_n \end{smallmatrix}]^T$ seems very unintuitive, but works for every holomorphic $\mathbb{R}^n \rightarrow \mathbb{R}^m$ function\footnote{for implementation this requires that the function can be evaluated with complex inputs} \cite{martins2021}.
Complex-step differentiation also is numerically stable, and has also an $\mathrm{O}(h^2)$ error size \cite{martins2003}. 
Both finite differences and complex-step scale linearly with the number of variables, but the finite difference method requires two function evaluations per variable, while complex-step requires only one.


\subsection{Numerical Integration methods}
To get a new discrete state from the continuous-time prediction process of the EKF, the nonlinear differential equations \ref{eq:predict_state} and \ref{eq:predict_cova} have to be solved.
As the last state estimate is available as an initial value, this can be expressed as the \textit{Initial Value Problem} (IVP)
\begin{align}
    \dot x(t) &= f(t, x, u) & x(0) &= x_{n} 
\end{align}
which is can be solved by numerical integration. 
The integration methods need to be explicit, as the filter in the prediction step has only past state information. 

Step size $\varepsilon$

\subsubsection{Explicit Runge-Kutta}

The explicit $s$-stage \textit{Runge-Kutta} method is compactly written as \cite{griffiths2010}
\begin{align}
    x_{n+1} &= x_n + \varepsilon_n \sum_{i=1}^{s} c_i k_i
    \\
    k_i &= f \left( t_n + a_i \varepsilon_n \, , \; x_n + \varepsilon_n \sum_{j=1}^{i-1} b_{ij} k_j \, , \ u \right) 
\end{align}
where $a_i$, $b_{ij}$, and $c_i$ are the entries of the vectors $a$ and $c$ and the matrix $B$, displayed in a \textit{Butcher tableau}
\begin{equation}
    \begin{array}{c|c}
        a & B \\
        \hline 
        & c^T
        \end{array}
    \label{eq:appendix_general_butcher}
\end{equation}
These can be very large, as in the \textit{Dormand-Prince} method (RK45, which the standard solver in Matlab).
For our cases we will keep the tableaus small.

The second order ($s=2$) butcher tableau of the Heun method \cite{griffiths2010}, which is used for the prediction covariance, is  
\begin{equation}
    \begin{array}{c|cc}
        0 &  &   \\
        1 & 1 &  \\
        \hline 
        & \frac{1}{2} & \frac{1}{2} 
        \end{array}
    \label{eq:appendix_heun_butcher}
\end{equation}

The fourth order ($s=4$) butcher tableau of the classical Runge-Kutta method RK4 \cite{griffiths2010}, which is used for solving the state prediction, is 
\begin{equation}
    \begin{array}{c|cccc}
        0 &  &  &  &  \\
        \frac{1}{2} & \frac{1}{2} &  &  &  \\
        \frac{1}{2} & 0 & \frac{1}{2} & & \\
        1 & 0 & 0 & 1 &  \\
        \hline 
        & \frac{1}{6} & \frac{1}{3} & \frac{1}{3} & \frac{1}{6}
    \end{array}
    \label{eq:appendix_rk4_butcher}
\end{equation}